\cdpchapter{Resumen}

BargAIn es una plataforma de compra inteligente orientada a optimizar la cesta de la compra
del consumidor combinando tres variables de decision: precio, distancia y tiempo.
El proyecto desarrolla una solucion full-stack con backend en Django/DRF,
base de datos PostgreSQL + PostGIS, cliente movil en React Native (Expo),
y companion web para el portal de comercios.

La propuesta integra cuatro capacidades diferenciales en una unica aplicacion:
(1) comparacion multi-tienda de precios,
(2) optimizacion multicriterio de ruta con OR-Tools,
(3) digitalizacion de listas y tickets mediante OCR con Google Cloud Vision,
y (4) asistente conversacional basado en modelo LLM via backend proxy.

Desde la perspectiva de ingenieria del software, el trabajo cubre el ciclo completo:
analisis de requisitos, diseno arquitectonico, implementacion modular,
pruebas unitarias/integracion/E2E, despliegue en staging y documentacion tecnica.
La arquitectura se organiza en dominios desacoplados
(users, products, stores, prices, shopping\_lists, optimizer, ocr, assistant, business,
notifications), con tareas asincronas en Celery y Redis para procesos pesados.

Como resultado, se obtiene un prototipo funcional validado en entorno de desarrollo y staging,
con cobertura de pruebas en modulos criticos, flujos E2E automatizados para web companion,
y verificaciones manuales en flujos moviles dependientes de GPS/camara.
La memoria documenta ademas limitaciones, decisiones arquitectonicas y lineas de evolucion
para convertir BargAIn en producto escalable.

\cdpchapter{Abstract}

BargAIn is an intelligent shopping platform focused on optimizing grocery baskets by combining
three decision variables: price, distance, and time. The project delivers a full-stack solution
with a Django/DRF backend, PostgreSQL + PostGIS geospatial storage,
React Native (Expo) mobile client, and a web companion for small businesses.

The proposal integrates four key capabilities in a single system:
(1) multi-store price comparison,
(2) multicriteria route optimization with OR-Tools,
(3) OCR-based list/ticket digitization using Google Cloud Vision,
and (4) a conversational assistant powered by an LLM through a backend proxy.

From a software engineering standpoint, the work covers the whole lifecycle:
requirements analysis, architecture design, modular implementation,
unit/integration/E2E testing, staging deployment, and technical documentation.

The final outcome is a functional prototype validated through automated and manual tests,
with clear traceability of technical decisions, current limitations,
and future evolution paths toward production readiness.


